\clearpage
\fancyhead[L]{\small GERLING.}
\hypertarget{native-547-lower}{}\label{native-547-lower}
\begin{center}\small Göttingische gelehrte Anzeigen. No. 190. pp. 1889..1891. November 28, 1812.\\[3pt]\rule{.72\linewidth}{.3pt}\end{center}

\emph{Methodi proiectionis orthographicae usum ad calculos parallacticos facilitandos explicavit, simulque eclipsin solarem die VII. Sept. \textup{1820} apparituram hoc modo tractatam mappaque geographica illustratam tamquam exemplum proposuit} \textsc{Chr. Lud. Gerling}. 46 pages in quarto. Göttingen. Printed by \textsc{I. C. Baier}.

To determine the circumstances of a stellar occultation or a solar eclipse for all points on Earth where it is visible, astronomers, as is well known, use orthographic projection onto a plane to which the straight line from the center of the occulted or eclipsed body is perpendicular. This procedure gives a clear, readily grasped picture of the different points in space involved and thus provides a convenient way to solve the problems that arise, whether one chooses drawing or calculation, the sharpest of all instruments. This short paper, which \SourceLeaf{548}its author --- formerly our learned fellow citizen and now a teacher of mathematics at the Royal Lyceum in Cassel --- submitted to the local Faculty of Philosophy as a qualifying dissertation for his doctorate, gives an accessible guide to the method and shows a thorough knowledge of the subject. The derivations are analytic almost throughout, so that the author found it unnecessary to include even a single figure. The fact that he could leave figures out without sacrificing clarity, and that anyone with even a little practice can readily follow his reasoning, demonstrates the clarity inherent in the analytical method when it is applied appropriately to geometrical problems.

The solar eclipse of September 7, 1820, to which the author applies the method as an example, will for a long time be one of the most remarkable in our part of Europe. The places where the eclipse is central lie on a line that begins in the northernmost part of America, passes to the north of Greenland and Iceland and through the North Sea, then crosses Germany in the direction from Emden to Trieste. It continues along the length of the Adriatic Sea and across Morea and the northernmost tip of the Red Sea, and ends at sunset in Arabia. Along this entire line and in a fairly broad zone above and below it, the eclipse is annular. This includes us here in Göttingen, where, according to the author's calculation, the ring will last from 2 hours 40 minutes 41 seconds to 2 hours 46 minutes 0 seconds, true time. The author also gives the main circumstances of the eclipse here for Bremen, Seeberg, Berlin, Vienna and Mannheim, and promises to give them elsewhere for a still larger number of locations.
